With the widespread adoption of AI applications, cyber-physical systems with autonomous capabilities, such as cars, drones and industrial robots, have become a prominent area of focus. In this fast-moving scene, being able to deploy autonomous systems on the latest and more capable hardware gains importance. Commercial Off-The-Shelf (COTS) systems provide good value for this purpose: they benefit from the performance improvements of the latest processor generations while staying accessible and cost-effective. For the implementation and rapid prototyping of autonomous systems, ROS~2 is a popular framework that is widely used in both industry and academia. Autonomous systems, however, must most of the time comply with deadline requirements to be considered functional, safe, or both; and because their target functionalities also require high-throughput computation, they frequently use hardware acceleration, such as GPUs.

Integrating GPUs into ROS~2 execution pipelines, however, introduces significant challenges. To promote rapid development and software reuse, ROS~2 provides an abstraction layer that effectively decouples applications from the underlying operating system~\cite{wu-2026-adssota}. In ROS~2, the unit of computation is a callback, activated by a timer or by the arrival of a message, a service request or response, or a waitable event. Callbacks are owned by run-time objects called nodes, and executors dispatch the ready callbacks of their assigned nodes on one or more threads of a process.

This deployment creates a case where executor threads are scheduled by the OS, ready callbacks are scheduled by executors, and, once a callback issues an API call to the GPU run-time, the resulting GPU work is scheduled by the GPU's internal scheduler, effectively creating two boundaries, between the OS and ROS~2 and between ROS~2 and the GPU run-time, across which no scheduler is aware of the others. Moreover, the order in which computations execute is determined by the logical structure of the application. This creates scenarios where CPU computations that depend on results from a GPU computation compete with each other for GPU resources and delay their downstream callbacks, which makes the analysis both hard and multi-layered.

Timing analysis of ROS~2 executors is a relatively new but developing area. The first analysis of the single-threaded executor was provided by Casini et al.~\cite{casini-19-reservationbased}, who modeled its polling-point/processing-window behavior and derived response-time bounds for processing chains under reservation-based OS scheduling. Successive works tightened these bounds with sharper interference characterizations~\cite{tang-2020-rtanalysispriorityassignmentchains}, execution-time curves~\cite{blass-2021-starvation} and DAG decompositions~\cite{tang-2023-rtaexecutor}. The analysis was later extended to the multi-threaded executor~\cite{jiang-22-processingchains,sobhani-2023-mtexecutoranalysis} and to reaction time and data age across multiple executors~\cite{teper-2024-end2endmtexecutoranalysisoptimization}. All of these analyses take the structure of the model as given: the chains, the callback-to-executor mapping, callback-group membership and priorities are inputs, and the parameters are either synthetic or measured with under-specified instrumentation; the taxonomy of a recent survey marks execution times as measured for only one of the listed analyses~\cite{casini-2025-surveyros2}. These analyses also model CPU execution only; as the survey puts it, \textit{``the vast majority of the discussed works focus on CPU scheduling only and neglect the presence of hardware accelerators''}~\cite{casini-2025-surveyros2}.

A second line of work extracts the model from a running system instead of assuming it. ROS-Llama~\cite{blass-2021-rosllama-paper} was the first to recover callback graphs and execution-time curves through tracepoints in rclcpp and rcl; Abaza et al.~\cite{abaza-2024-ros2modelsynthesis} synthesize callback DAGs with eBPF probes on the ROS~2 libraries, assuming single-threaded executors; LiME~\cite{brandenburg-25-lime} and LiMEIT~\cite{brandenburg-2026-limeit} take a framework-agnostic route, inferring thread-level and, most recently, intra-thread task models from scheduler events and sampled call stacks, again for single-threaded executors only. All three measure execution as consumed processor time, and none observes the GPU layer or multi-threaded executors. Tracing infrastructure specialized for ROS~2 also exists: ros2\_tracing~\cite{bedard-2022-ros2tracing} instruments rclcpp, rcl and rmw with low-overhead LTTng tracepoints, and extensions such as Autoware\_Perf~\cite{li-2022-autowareperf} add inter-node latency on top of it but rely on hand-written architecture files.

On the GPU side, experimental work on streams and queues~\cite{otterness-17-inferringpolicies,amert-17-tx2review} and on channels, runlists and engines~\cite{bakita-24-demystify} turned the GPU scheduler from a black box into a documented hierarchy, on which our GPU-side model builds. A prominent approach for managing the GPU under ROS~2 is to put an arbitration layer between ROS~2 and the GPU, such as PAAM~\cite{enright-2024-paam} and ROSGM~\cite{li-2023-rosgm}, which moves accelerator calls out of the callback into a server that schedules them by chain priority. The server must be told which callbacks use which accelerator and for how long, which is exactly what FISH recovers from traces. No existing approach observes the OS, executor and GPU layers of a deployed ROS~2 system together.

To provide an end-to-end analysis for an autonomous ROS~2 system running on a COTS platform, therefore, a model synthesis method is required that extracts the application's logical structure together with the execution characteristics of its callbacks and the scheduling dependencies across layers. FISH aims to provide such a holistic model of the system: every callback is mapped to the higher-level entities that own it, nodes and executors, together with the interdependencies between callbacks; and where a callback uses the GPU, the model reveals the execution sequence of its CPU segments, copy-engine transfers and kernel executions on the compute engine, with the dependencies between them. We make the following contributions: (I)~a hierarchical model of a deployed ROS~2 system spanning the OS, ROS~2 and GPU layers; (II)~instrumentation that extends ros2\_tracing with tracepoints for object lifetime, all five callback dispatch kinds and intra-process delivery, covering both rclcpp and rclpy; and (III)~a trace-acquisition method for containerized deployments: launch-time wrapping of GPU processes, a system-stable signal for reproducible runs, and phase boundaries (initialization, warm-up, steady state) derived from the trace itself.
